\documentclass[10pt,a4paper]{amsart}
\usepackage[margin=19mm]{geometry}
\usepackage{amsmath,amssymb,mathtools}
\usepackage[scr=rsfs,cal=euler]{mathalfa}
\usepackage{xfrac}
\usepackage[unicode,hidelinks,bookmarks=false]{hyperref}
\newcommand{\ie}{i.e.\ }
\newcommand{\cf}{cf.\ }
\newcommand{\resp}{respectively\ }
\newcommand{\Subsection}{\subsection}
\providecommand{\nobreakdash}{\nobreak\hbox{-}}
\setlength{\emergencystretch}{2em}
\multlinegap0pt
\usepackage{amssymb}
\newtheorem{definition}{Definition}
\newtheorem{thm}{Theorem}
\newtheorem{cor}{Corollary}
\def\Cl{\mathop{\rm Cl}\nolimits}
\def\Div{\mathop{\rm Div}\nolimits}
\def\Q{{\mathbb Q}}
\def\supp{\mathop{\rm Supp}\nolimits}
\title{On the Boucksom-Zariski decomposition for irreducible symplectic varieties and bounded negativity}
\author{Micha{\l} Kapustka, Giovanni Mongardi, Gianluca Pacienza, Piotr Pokora}
\date{Source: arXiv:1911.03367v6 (CC BY 4.0)}
\begin{document}
\maketitle

\noindent\emph{Source and licence.} On the Boucksom-Zariski decomposition for irreducible symplectic varieties and bounded negativity. Version arXiv:1911.03367v6 (CC BY 4.0), distributed by arXiv under the Creative Commons Attribution 4.0 International licence, http://creativecommons.org/licenses/by/4.0/, as recorded in the arXiv licence metadata for this exact version and shown on the abstract page https://arxiv.org/abs/1911.03367v6. Full licence notice: \texttt{licenses/arxiv-1911.03367-CC-BY-4.0.txt}.\par\smallskip

\noindent\textit{Editorial scope.} Counted unit: the article's cor cor (source0.tex lines 649--650). The article's complete original proof of it is delivered with it (source0.tex lines 651--653, one \texttt{proof} environment of 3 lines). No mathematics is written, repaired or re-derived: the statement and the proof are the article's own lines, in the article's own notation, and no retained line is substituted or re-flowed.  Retained uncounted as the source-local context the proof needs: lines 563--574; lines 598--601; lines 627--629.  Nothing was omitted from any retained span, so the package carries no omission record.  No retained line contains a citation, so the wrapper carries no bibliography and no bibliography entry.  No retained line contains a figure environment, an image-inclusion command or drawing code, so no image is delivered and the package declares no asset dependency.  Editorial formatting only, in addition: an AMS \texttt{amsart} preamble that restates the article's own declarations and macros used by the retained lines, namely the environments \texttt{definition}, \texttt{thm}, \texttt{cor} on separate counters, together with the standard packages those lines need.  The retained environments print in this document as Definition 1, Theorem 1, Corollary 1 in order of appearance, whereas the article prints the same items with the numbering of its own full document.  The complete native source of the article as distributed in the arXiv e-print for this exact version is bundled unchanged as \texttt{sources/102982/source0.tex}.
\begin{definition} \label{prime} 
Let $X$ be a normal compact complex space endowed with a quadratic form $q_{X}$ on $\Div_{\Q}(X)=\Div(X)\otimes \Q$, where $\Div(X)$ denotes the group of Weil divisors (without any equivalence relation).

A reduced and irreducible effective divisor $D\subset X$ is called a prime divisor. An effective Weil $\Q$-divisor $E=\sum a_i E_i$ 
is called $q_{X}$-exceptional if the Gram matrix   $(q_{X}(E_i,E_j))_{i,j}$ of the irreducible components of the support of $E$ 
is negative definite.  
Furthermore, we say that a Weil $\Q$-divisor $D$ is $q_{X}$-nef if 
$$q_{X}(D,E)\geq 0 $$ 
for every effective 
Weil $\Q$-divisor $E$, or equivalently for every prime divisor $E$.
In this context, we will say that $q_X$ is an intersection product (without mentioning any basis) if it is an intersection product with respect to the basis consisting of all prime divisors.
\end{definition}

\begin{thm}\label{thm:q-Zar}
Let $X$ be a normal compact K\"ahler variety endowed with a quadratic form $q_{X}$ on  $\Div_{\Q}(X)$ which is an intersection product  (in the sense of Definition \ref{prime}).
Then all effective Weil $\Q$-divisors on $X$ have a unique rational  $q_{X}$-Zariski decomposition. Moreover, if $q_X$ is compatible with linear equivalence then it also induces a unique $q_X$-Zariski decomposition on $\Cl_{\Q}(X)$ and both decompositions are compatible.
\end{thm}

\begin{equation}\label{eq:negative_support}
 N_i \subset \supp(M).
\end{equation}

\begin{cor} Let $q_X$ be an intersection product on $\Div_{\Q}(X)$ compatible with linear equivalence.  The decomposition  $D= P(D) + N(D)$ provided by Theorem \ref{thm:q-Zar} for an effective Weil $\Q$-divisor $D$  satisfies $\mathrm{Supp} (P(D))\cup \mathrm{Supp} (N(D))=\mathrm{Supp} (D)$. In particular, for the $q_X$-Zariski decomposition $[D]=[P]+[N] $ in $\Cl_{\Q}(X)$, we can choose $P(D)\sim_{\mathrm{lin}} P$ and  $N(D)\sim_{\mathrm{lin}} N$ such that $\mathrm{Supp} (P(D))\cup \mathrm{Supp} (N(D))=\mathrm{Supp} (D)$.
\end{cor}

\begin{proof}
The inclusion $\mathrm{Supp} (P(D))\subset \mathrm{Supp}(D)$ is given by the natural map of sections, while the inclusion $\mathrm{Supp} (N(D))\subset \mathrm{Supp}(D)$ is given by \eqref{eq:negative_support}. The other inclusion follows from the fact that both $P(D)$ and $N(D)$ are effective.
\end{proof}

\end{document}
